\section{Preliminaries: Continuous Flow Language Modeling and RL}
\label{sec:estimators}

\paragraph{Continuous flow language modeling.}
We study ELF \citep{hu2026elf}, a continuous flow language model that generates text in embedding space. A token sequence $w=(w_1,\dots,w_L)$ is embedded by a frozen T5 encoder \citep{raffel2020t5} as $x=\mathrm{Enc}(w)\in\mathbb R^{L\times d}$, one vector $x_i\in\mathbb R^d$ per position. The first positions hold the prompt $c$ and stay clean throughout; the remaining positions, the set $\mathcal G$, are generated. ELF is trained along the linear interpolation $z_t=(1-t)c_\sigma\varepsilon+tx$, where
$\varepsilon\sim\mathcal N(0,I)$, $t\in[0,1]$, and $c_\sigma$ is a fixed noise scale. Its denoiser transformer reads $z_t$ and $t$, with the prompt in the prefix positions, and outputs a new estimate $\hat x$:
\begin{equation}
\hat x=\hat X_\theta(z_t,t)\in\mathbb R^{L\times d}; \quad v_\theta(z_t,t)=\frac{\hat x-z_t}{1-t}.
\label{eq:denoiser}
\end{equation}
We write $\hat X_\theta$ for one denoiser call. At the endpoint, a decoder maps each embedding to vocabulary logits $\ell_i\in\mathbb R^{|\mathcal V|}$, yielding $w_i=\arg\max\ell_i$. 
% Throughout, $\mathrm{sig}(\cdot)$ denotes the logistic function and $\mathrm{logit}=\mathrm{sig}^{-1}$.

\paragraph{The SDE sampler.}
ELF provides a stochastic sampler in which, given $z_{k-1}$, each step perturbs the state with noise controlled by the churn parameter 
$\gamma$ and then integrates the velocity forward:
\begin{equation}
y_k=\alpha_k z_{k-1}+\sigma_k\varepsilon_k, z_0:=y_0, \hat x_k=\hat X_\theta(y_k,\tilde t_k), z_k=y_k+\lambda_k\big(\hat x_k-y_k\big)=:D_\theta(y_k)
\end{equation}
where $\tilde t_k=\alpha_k t_{k-1}$, $\alpha_k=1-\gamma (t_k-t_{k-1})$, $\sigma_k=(1-\alpha_k)c_\sigma$, $\lambda_k=\frac{t_k-\tilde t_k}{1-\tilde t_k}$ and $\varepsilon_k \sim \mathcal N(0,I)$. The final step is deterministic followed by the decoder. The states $y_k$ form a Markov chain with trajectory $\tau=(y_k)_{k=1}^{K-1}$ and exact Gaussian density:
\begin{equation}
y_{k,i}=\alpha_kD_\theta(y_{k-1})_i+\sigma_k\varepsilon_{k,i}, \quad p_\theta(\tau)=p(y_1)\prod_{k=2}^{K-1}\prod_{i=1}^{|\mathcal G|}\mathcal N\!\left(y_{k,i};\alpha_kD_\theta(y_{k-1})_i,\sigma_k^2I_d\right)
\label{eq:chain}
\end{equation}
where $k$ and $i$ denote the flow step and token position and $y_1=\alpha_1y_0+\sigma_1\varepsilon_1$ is Gaussian noise whose law $p(y_1)$ does not depend on $\theta$. The response $w=w_\theta(y_{K-1})$ is obtained deterministically by advancing $z_{K-1}=D_{\theta}(y_{K-1})$ to $t_K=1$ followed by the decoder.

\paragraph{RL objective and value.}
A program verifier assigns a decoded response $w$ a binary reward $r(w)$. RL post-training maximises the expected reward,
\begin{equation}
J(\theta)
=
\mathbb E_{\tau\sim p_\theta(\cdot)}
\left[
r\bigl(w_\theta(\tau)\bigr)
\right].
\label{eq:objective}
\end{equation}
Because subsequent flow steps inject Gaussian noise, the same intermediate state $y_k$ can produce different responses. Thus the verifier outcome conditioned on $y_k$ is random:
\begin{equation}
r\bigl(w_\theta(\tau)\bigr) \mid y_k, t_k
\sim
\operatorname{Bernoulli}\!\left(V_k^*(y_k)\right),
\qquad
V_k^*(y_k)
=
\mathbb E_{\tau\sim p_\theta(\cdot\mid y_k,t_k)}
\!\left[r\bigl(w_\theta(\tau)\bigr)\right].
\label{eq:value}
\end{equation}
The terminal value $V_{K-1}^*(y_{K-1})=r(w_\theta(y_{K-1}))$ is piecewise constant, whereas earlier values depend on $\theta$ through subsequent transitions and satisfy $V_{k-1}^*(y_{k-1})=\mathbb E_{\varepsilon_k} \left[V_k^*(y_k)\mid y_{k-1}\right]$. Gaussian re-noising smooths the bounded terminal reward over intermediate states (Appendix~\ref{app:pw_derivation}).

\paragraph{Two gradient estimators.}
The \emph{score-function} estimator differentiates the density: $\nabla_\theta J=\mathbb E_\tau[A\,\nabla_\theta\log p_\theta(\tau)]$, where $A(\tau)=r(w_\theta(\tau))-b$ and any baseline $b$ independent of the sampled trajectory leaves the expected gradient unchanged \citep{williams1992reinforce}. For an AR model, $\log p_\theta(w)=\sum_j\log\pi^{\rm AR}_\theta(w_j\mid w_{<j})$, so each term corresponds to a categorical token decision. For the flow chain~(\ref{eq:chain}), the same estimator applies to the exact Gaussian density. The \emph{pathwise} estimator instead differentiates the sample. Writing the sampler as $\tau=F_\theta(\varepsilon)$, the pathwise estimator differentiates a smooth reward estimate $g$: $\nabla_\theta\mathbb E[g(\tau)]=\mathbb E_\varepsilon[\nabla_\theta g(F_\theta(\varepsilon))]$ under regularity conditions. It needs a differentiable path from $\theta$ to $g$. The terminal reward is nondifferentiable due to decoder $\arg\max$ and program verification, whereas the state value in~(\ref{eq:value}) is differentiable.
